\documentclass[xcolor=dvipsnames]{beamer}

% 1. PREÁMBULO: Configuraciones Globales
\usetheme{Madrid}
\usepackage[utf8]{inputenc}
\usepackage[spanish, es-tabla]{babel} % es-tabla para que diga Tabla y no Cuadro
\usepackage{amsmath, amssymb}
\usepackage{diagbox}
\usepackage{multirow}
\usepackage{tikz}

% 2. CONFIGURACIÓN DE TIKZ (Debe ir en el preámbulo)
\usetikzlibrary{arrows.meta, positioning, calc}
\tikzset{
    state/.style={circle, draw=black, thick, minimum size=1.1cm, font=\bfseries},
    alc/.style={state, fill=Emerald!20, text=Emerald!70!black},
    baj/.style={state, fill=Maroon!20, text=Maroon!70!black},
    est/.style={state, fill=Gray!20, text=Gray!70!black},
    trans/.style={->, >={Stealth[scale=1.0]}, thick, black!70, shorten >=1pt},
    val/.style={font=\tiny\bfseries, text=blue!80!black, yshift=2mm}
}

% 3. DATOS DE LA PRESENTACIÓN
\title{Procesos de Decisión de Markov}
\author{Daniel}
\date{2026}

% 4. CUERPO DEL DOCUMENTO
\begin{document}

\maketitle

\begin{frame}{Cálculo de la Función Valor: Recursión}
Para calcular $V_t(s)$, propagamos los valores desde el horizonte final ($t=0$) hacia atrás.

\begin{center}
\begin{tikzpicture}[node distance=1.5cm, scale=0.9]
    % 1. Nodos de tiempo (Siempre visibles o desde el inicio)
    \node[font=\small] at (0, 2.5) {$t=2$};
    \node[font=\small] at (3, 2.5) {$t=1$};
    \node[font=\small] at (6, 2.5) {$t=0$ (Fin)};
    
    % 2. Columna t=0 (Caso Base) - Aparece en la transparencia 2
    \node<2->[alc] (A2) at (6, 1.5) {A};
    \node<2->[baj] (B2) at (6, 0) {B};
    \node<2->[val] at (A2.north) {$V_0=0.005$};
    \node<2->[val] at (B2.north) {$V_0=-0.005$};

    % 3. Columna t=1 - Aparece en la transparencia 3
    \node<3->[alc, fill=white, draw=black!30] (A1) at (3, 1.5) {A};
    \node<3->[baj, fill=white, draw=black!30] (B1) at (3, 0) {B};

    % 4. Flechas de cálculo V1 - ¡CRÍTICO! Deben ser <3-> para que A1 ya exista
    \draw<3->[trans, blue] (A1) -- (A2) node[midway, above, font=\tiny] {0.9};
    \draw<3->[trans, blue!50] (A1) -- (B2) node[pos=0.7, above, font=\tiny] {0.1};
    
    % Valor calculado para V1 - Aparece en la transparencia 4
    \node<4->[val] at (A1.north) {$V_1=0.009$};

    % 5. Columna t=2 - Aparece en la transparencia 5
    \node<5->[alc] (A0) at (0, 1.5) {A};
    \node<5->[baj] (B0) at (0, 0) {B};
    
    % Flecha que conecta t=2 con t=1
    \draw<5->[trans, red] (A0) -- (A1) node[midway, above, font=\tiny] {0.9};
\end{tikzpicture}
\end{center}

\begin{itemize}
    \item<2-> \textbf{Paso 0:} $V_0(s) = r(s)$ (Caso base).
    \item<3-> \textbf{Paso 1:} Calculamos $V_1$ usando los nodos de la derecha.
    \item<5-> \textbf{Paso 2:} Repetimos hacia atrás para $V_2, V_3 \dots$
\end{itemize}
\end{frame}

\end{document}